\chapter{信息感知与状态估计}
\label{ch:est}

被动观测对内部热阻的信息在什么条件下为零，决定探针是否必要。规格中的估计器是分层集合滤波；本次日志里的估计器是一个标量跟踪器。跟踪器只看见表面温度和平均含水，时间常数正比于热阻估计和厚度的平方。第一个记录点出现负值，来自先验热阻 5 的滞后上限：估计从冷点之下爬起，经过探针窗之后才越过零。这一滞后是开局余量，写入附录轨迹表。

\section{可安装观测}

进入品质估计器的通道列于表\ref{tab:obs}。安全侧测量不进入品质滤波。腔温、湿度、质量和红外都可以作为可安装通道写入观测向量；工程核实际读入的是表面温度加 \(0.4^\circ\mathrm{C}\) 散列，以及平均含水加 \(0.004\) 散列。其余通道留在表里，用来标明“可安装”与“已经接到跟踪器”的差别。

\begin{table}[htbp]
\centering
\caption{进入品质估计器的观测}
\label{tab:obs}
\begin{tabularx}{\textwidth}{llX}
\toprule
通道 & 符号 & 与隐藏状态的关系 \\
\midrule
腔体空气温度 & \(y_T\) & 靠近 \(T_\infty\)，受加热器辐射污染 \\
腔体相对湿度 & \(y_\varphi\) & \(p_v/p_{\mathrm{sat}}(T_\infty)\)，饱和时丢失分辨率 \\
食品质量 & \(y_m\) & 累计蒸发的积分，不定位冷点 \\
表面热像 & \(\vect{y}_{\mathrm{IR}}\) & 发射率未知的辐射亮度 \\
电功率 & \(y_P\) & 已知输入，写入边界 \\
排气温度 & \(y_{T,e}\) & 冷凝与漏汽的积分证据 \\
\bottomrule
\end{tabularx}
\end{table}

观测方程 \(\vect{y}_k=h\big(\Phi(\vect{X}(t_k)),\vect{\theta}\big)+\vect{v}_k\)。红外通道
\[
L_j=\epsilon(\varepsilon_{l,s},S_s)\,\sigma T(\vect{r}_j)^4
+\big(1-\epsilon(\varepsilon_{l,s},S_s)\big)L_j^{\mathrm{refl}}.
\]
发射率随表面变干下降。系统精度应取发射率区间两端分别滤波后的分位包络。本次核没有热像，也没有质量。进入跟踪器的是表面温度加 \(0.4^\circ\mathrm{C}\) 的确定性散列，以及平均含水加 \(0.004\) 的散列。噪声由种子和步号生成，重复运行得到同一轨迹。

\section{病态结构}

线性化热方程在厚度方向上，对高频边界噪声的放大随 \(\mathrm{e}^{kL}\) 增长\cite{ozisik2000}。半无限体上表面温度取成角频率 \(\omega\) 的谐波时，
\begin{equation}
\label{eq:halfspace}
T(x,t)=\mathrm{e}^{-kx}\cos(\omega t-kx),\qquad
k=\sqrt{\frac{\omega}{2\alpha}}.
\end{equation}
深度 \(L\) 处的振幅是表面振幅的 \(\mathrm{e}^{-kL}\)。\(\alpha=2.0\times 10^{-7}\,\mathrm{m}^2/\mathrm{s}\)。表\ref{tab:amp}给出正文中已经出现的时间尺度。

\begin{table}[htbp]
\centering
\caption{半无限体上单个谐波的放大倍数 \(\mathrm{e}^{kL}\)}
\label{tab:amp}
\begin{tabular}{llrr}
\toprule
周期 & 位置 & \(L=30\,\mathrm{mm}\) 放大 & 穿透深度 \\
\midrule
8 s & 演示记录间隔 & \(1.81\times 10^{18}\) & 0.71 mm \\
40 s & 探针窗 & \(1.46\times 10^{8}\) & 1.60 mm \\
1170 s & 30 mm 开环菜单 & 32.3 & 8.63 mm \\
\bottomrule
\end{tabular}
\end{table}

8 s 谐波在 15 mm 处剩下 \(7.44\times 10^{-10}\,\mathrm{K}\)。菜单时间把放大降到另一个量级，薄件在自己的菜单时间上还留着约十分之一，厚件只留到千分之几。令 \(kL=1\) 对应的周期是 \(\pi L^2/\alpha\)：30 mm 为 14137 s，长于菜单。加密表面像素提高的是已经衰减掉的高频分量；要看见冷点，必须有一个穿透时间尺度上的动态激励，或一个对内部热阻敏感的积分量。质量损失对累计蒸发敏感，对冷点位置不敏感。图\ref{fig:probe5}给出种子 7 探针窗前后的表面、冷点与估计。窗内冷点几乎不动，估计仍为负值。

\begin{figure}[htbp]
\centering
\includegraphics[width=0.84\textwidth]{P_probe.png}
\caption{种子 7 探针窗口。冷点从 \(8.01^\circ\mathrm{C}\) 到 \(8.05^\circ\mathrm{C}\)；估计在窗结束时刚越过零。}
\label{fig:probe5}
\end{figure}

图\ref{fig:d23sd}给出 15 mm 同一格子上 30 个种子的停止时刻。开环钉在菜单 900 s；腔温比例钉在 3300 s；无探针集中在 1050 s 附近；对偶多数与无探针重合，少数种子提前到 650 s 量级。提前来自跟踪器先于真值进度越过阈值：15 mm 的 Fourier 数在菜单时刻已经接近 1，内部热阻对停机时刻的贡献被表面进度饱和盖住。

\begin{figure}[htbp]
\centering
\includegraphics[width=0.84\textwidth]{D23_seed_stop.png}
\caption{15 mm 同一格子、30 个种子的停止时刻。第二列为腔温比例，第三列为无探针。}
\label{fig:d23sd}
\end{figure}

\section{稳态不可辨识}

\begin{proposition}[平衡点不含内部参数]
\label{prop:steady}
设\eqref{eq:twonode}中 \(\dot m=0\)、\(q_{\mathrm{lat}}=0\)，且存在平衡点。则 \(T_1=T_2=T_\infty\)。\((R_{\mathrm{int}},C_1,C_2)\) 不进入该代数方程，平衡观测对三者的灵敏度为零。
\end{proposition}

\begin{proof}
第二式在平衡时为 \((T_1-T_2)/R_{\mathrm{int}}=0\)。热阻取正，故 \(T_1=T_2\)。第一式变为 \((T_\infty-T_1)/R_{\mathrm{ext}}=0\)，故 \(T_1=T_\infty\)。三个参数只作为除数出现，不出现在解的表达式中。对任一参数求导，平衡方程的右端仍为零。
\end{proof}

推论：恒温段再密的腔温采样也不更新内部热阻。信息动作必须离开平衡。断言对 \(R\in\{1,3,6\}\) 从同一初值积分，平衡温差的极差为 0。这是“到时按菜单停机”失败的数学形式。

灵敏度 \(\vect{\sigma}=\partial(T_1,T_2)/\partial R_{\mathrm{int}}\) 满足
\begin{equation}
\label{eq:sigma}
\dot{\vect{\sigma}}=\vect{A}\vect{\sigma}+\frac{\partial\vect{A}}{\partial\theta}\begin{pmatrix}T_1\\ T_2\end{pmatrix}.
\end{equation}
平衡且 \(T_1=T_2\) 时右端为零，\(\vect{\sigma}=0\) 是平衡解，Fisher 矩阵里这一列贡献为零。离开平衡、温差不为零时，右端不再是零。

\section{Fisher 信息与在线跟踪器}

参数灵敏度满足
\begin{equation}
\label{eq:sens}
\frac{\mathrm{d}}{\mathrm{d}t}\frac{\partial \vect{s}}{\partial\theta_i}
=\vect{A}(\vect{\theta})\frac{\partial \vect{s}}{\partial\theta_i}
+\frac{\partial\vect{A}}{\partial\theta_i}\vect{s}
+\frac{\partial\vect{B}}{\partial\theta_i}\vect{u}.
\end{equation}
Fisher 信息
\begin{equation}
\label{eq:fisher}
\vect{\mathcal{I}}(\vect{\theta};\vect{u}_{0:K})
=\sum_{k=0}^{K}
\Big(\frac{\partial h}{\partial \vect{s}}\frac{\partial \vect{s}_k}{\partial\vect{\theta}}\Big)^{\!\top}
\vect{R}_k^{-1}
\Big(\frac{\partial h}{\partial \vect{s}}\frac{\partial \vect{s}_k}{\partial\vect{\theta}}\Big).
\end{equation}
探针设计在规格上最大化 \(\log\det\vect{\mathcal{I}}\) 的增量。本次没有组装这个矩阵。

在线一侧读表面温度和平均含水。内部时间尺度
\[
\tau=\mathrm{clip}(\hat r,0.8,8)\,(L/0.03)^2\,280\ \mathrm{s}.
\]
估计冷点
\begin{align}
\hat T_c &\leftarrow \hat T_c+\Delta t\,(y_T-\hat T_c)/\tau,\\
\hat T_c &\leftarrow \min\big(\hat T_c,\, y_T-4\max(\hat r-1,0.2)\big).
\end{align}
第二式把估计压到表面读数之下。先验 \(\hat r=5\)，方差 \(0.64\)。种子 7 第一个记录点的估计是 \(-7.92^\circ\mathrm{C}\)：表面读数减去 \(4\times(5-1)=16\,\mathrm{K}\) 得到负值，开局时表面还带着很大的滞后余量。1200 s 时估计 \(61.27^\circ\mathrm{C}\) 高于冷点 \(46.65^\circ\mathrm{C}\)，滞后上限用表面减一项只依赖 \(\hat r\) 的常数，表面已经远高于冷点之后，估计仍可以停在冷点之上。已实现的跟踪器给出带滞后的点估计。交叉表里中等厚度对偶的中位误差接近零、厚件仍然为负，要放在这个观测器的偏差结构上看。

跟踪器每个时间步只更新一个标量。方差只在探针窗结束、且新热阻更接近真值参照时乘一次 0.35。集合卡尔曼滤波的分析步、观测 Gram 和 Fisher 矩阵写在估计器的完整形式里；主日志使用的是这一条标量规则。下游策略把 \(\hat T_c\) 放进饱和线性律，接口比九维决策状态窄，交叉表比较的就是这个跟踪器给出的停机偏差。

“看不见”在这一章落到两条可以核对的事实上。平衡点上三个温度对 \(R_{\mathrm{int}}\) 的导数为零，恒温段再密的腔温采样也不更新内部热阻。半无限体放大表给出另一条：8 s 谐波在 30 mm 处放大 \(1.81\times 10^{18}\)，40 s 探针窗仍有 \(1.46\times 10^{8}\)，开环菜单时间量级才降到几十；像素再密，提高的仍是已经衰减掉的高频分量。质量损失对累计蒸发敏感、对冷点位置不敏感，S2 把它写成零效应。两条事实指向同一动作：离开平衡、短于内部时间常数的两段脉冲。种子 7 在窗内冷点只升 \(0.04\,\mathrm{K}\)，估计在窗结束时刚越过零，1200 s 时估计 \(61.27^\circ\mathrm{C}\) 高于冷点 \(46.65^\circ\mathrm{C}\)，方差在窗后由 0.64 乘到 0.224。点估计可以先于真值进度宣布到达，厚件的导热时间尺度却走不完——交叉表里 30 mm 对偶最接近零、50 mm 仍为负，要放在这个观测器的偏差结构上看。规格层的分层集合滤波留在第\ref{ch:conc}章；交叉表读取的是日志层。
